%% ch16 --- Chaos wewnątrz komputera
%%
%% Napisany we wrześniu 2026. Każdy obraz chaosu w tej książce narysowała
%% maszyna, która zaokrągla w każdym kroku; ten rozdział pyta, co to robi, i
%% to mierzy. Liczby zmiennoprzecinkowe od zera; odwzorowanie podwajające
%% ginące w komputerze; jedna reguła zapisana na dwa sposoby, rozchodząca się
%% w ciągu mniej więcej pięćdziesięciu kroków, wobec prawdy liczonej ze stu
%% cyframi; dlaczego arytmetyka dokładna i rozszerzona nie ratuje długiego
%% chaotycznego przebiegu; ratunek, który działa (cieniowanie, sformułowane w
%% ramce i pokazane przez znalezienie cienia i przez zgodne statystyki);
%% wreszcie skończone maszyny, pętle i generatory liczb losowych. Każda
%% liczba pochodzi z code/measure/ch16.py, który importuje te same funkcje,
%% które drukują listingi.
%%
%% Bliźniak: chapters/en/ch16-computers.tex. Podrozdział za podrozdziałem,
%% makro za makrem, liczba za liczbą; tools/parity.py je porównuje.

\chapter{Chaos wewnątrz komputera}
\label{ch:computers}
\index{liczby zmiennoprzecinkowe}\index{zaokrąglanie}

\noindent
Każdy obraz chaosu w tej książce narysował komputer i tak samo policzył
każdą liczbę. Oto niewygodny fakt o maszynie, która je narysowała. Nie umie
zapamiętać liczby $\num{0.1}$. Zapamiętuje liczbę bardzo do niej zbliżoną, a
za każdym razem, gdy dodaje albo mnoży dwie liczby, zaokrągla wynik do
takiej, którą umie zapamiętać.

Przez kilka rozdziałów uczyłeś się, że chaos zamienia najmniejszy błąd w
duży. Skoro więc komputer robi drobny błąd w każdym kroku, jak którykolwiek
z tych obrazów może być poprawny? Ten rozdział mierzy odpowiedź, a jest ona
lepsza, niż mógłbyś się obawiać, i gorsza, niż mógłbyś mieć nadzieję.

\begin{outcomes}
\outcome{powiedzieć, co zapamiętuje komputer, gdy wpisujesz liczbę, i
  dlaczego $\num{0.1} + \num{0.2}$ nie daje $\num{0.3}$}
\outcome{przewidzieć krok, w którym odwzorowanie podwajające ginie w
  komputerze}
\outcome{oszacować, przez ile kroków chaotyczne obliczenie pozostaje
  poprawne, na podstawie liczby cyfr dwójkowych, które zachowuje}
\outcome{powiedzieć, którym wynikom chaotycznej symulacji możesz ufać, i
  dlaczego}
\outcome{opisać generator liczb losowych jako regułę deterministyczną i
  powiedzieć, dlaczego musi się powtarzać}
\end{outcomes}

\section{Co zapamiętuje komputer}
\label{sec:ch16-keeps}
\index{system dwójkowy}\index{podwójna precyzja}

\noindent
Zapisz jedną trzecią jako ułamek dziesiętny, a nigdy nie skończysz:
$\num{0.333}\ldots$, trójki ciągną się bez końca. Ułamki dziesiętne zapisują
dokładnie tylko te ułamki, których mianownik składa się z dwójek i piątek,
czynników dziesiątki. Komputer liczy zamiast tego dwójkami. W \emph{systemie
dwójkowym} miejsca po przecinku są warte połowę, ćwierć, jedną ósmą i tak
dalej, więc $\tfrac{3}{4}$ zapisuje się jako $\num{0.11}$: połowa plus ćwierć.

\begin{think}
Które z tych liczb system dwójkowy zapisze dokładnie, stałą liczbą miejsc:
$\tfrac{1}{2}$, $\tfrac{3}{8}$, $\tfrac{1}{10}$?
\end{think}
\reveal{$\tfrac{1}{2}$ i $\tfrac{3}{8}$, ale nie $\tfrac{1}{10}$. System
dwójkowy zapisuje dokładnie ułamki, których mianownik jest potęgą dwójki, a
dziesięć ma czynnik pięć.}
Dwójkowo jedna dziesiąta to $\num{0.000110011}\ldots$, a blok $0011$
powtarza się bez końca, więc komputer musi ją gdzieś uciąć. Python, jak
prawie każdy program, którego używasz, przechowuje liczbę jako
\emph{double}, liczbę podwójnej precyzji: \val{ch16.bits.double} cyfry
dwójkowe znaczące, czyli około \val{ch16.decimal.double} dziesiętnych,
plus potęgę dwójki, która mówi, gdzie stoi przecinek. Dla jednej dziesiątej
zapamiętuje najbliższą liczbę podwójnej precyzji, za dużą o mniej więcej
\val{ch16.tenth.err}. Każda suma i każdy iloczyn są potem także
\emph{zaokrąglane}: dokładny wynik zastępuje najbliższa liczba podwójnej
precyzji.

\begin{think}
Zarówno $\num{0.1}$, jak i $\num{0.2}$ są zapamiętane odrobinę błędnie, a ich
suma jest jeszcze raz zaokrąglana. Czy komputer zgodzi się, że
$\num{0.1} + \num{0.2}$ to $\num{0.3}$?
\end{think}
\reveal{Nie. Suma drukuje się jako liczba o włos większa od $\num{0.3}$, a
porównanie odpowiada \code{False}.}

\pyregion{code/ch16/floats.py}{store}{Co zapamiętuje komputer i co z tym robi}{lst:ch16-floats}

\transcript{ch16-floats}

\noindent
Trzeci wiersz to dokładna wartość zapamiętana dla $\num{0.1}$; czwarty
zapisuje ją jako ułamek, którego mianownik to $2^{\val{ch16.tenth.power}}$.
Każda liczba podwójnej precyzji to liczba całkowita podzielona przez potęgę
dwójki. Szósty wiersz dodaje trzy liczby w dwóch kolejnościach i dostaje dwa
wyniki: gdy każdy krok zaokrągla, grupowanie, o którym szkoła mówiła, że nie
ma znaczenia, ma znaczenie. Ostatni pokazuje, że brakuje miejsca: powyżej
$2^{53}$ liczba podwójnej precyzji nie mieści już każdej liczby całkowitej.

\begin{note}
Te reguły zaokrąglania ustalono w 1985 roku w standardzie IEEE~754, którego
przestrzega prawie każdy komputer. Wymaga on, by $+$, $-$, $\times$, $\div$ i
pierwiastek kwadratowy dawały dokładny wynik zaokrąglony do najbliższej
liczby podwójnej precyzji. Dlatego twój komputer powinien się zgadzać z każdą
liczbą wydrukowaną w tej książce: jej własne pomiary używają tylko tych
działań.
\end{note}

\section{Odwzorowanie, którego komputer nie umie policzyć}
\label{sec:ch16-doubling}
\index{odwzorowanie podwajające}

\noindent
Najprostsza chaotyczna reguła, jaka istnieje, mieści się w jednym wierszu:
podwój liczbę, a jeśli wynik wynosi $1$ albo więcej, odejmij $1$.
Symbolicznie $x \to 2x \bmod 1$, czyli \enquote{$2x$ bez części całkowitej}.
To \emph{odwzorowanie podwajające} w każdym kroku podwaja odstęp między
dowolnymi dwoma bliskimi punktami startowymi. W systemie dwójkowym widać, jak
działa: podwojenie przesuwa każdą cyfrę o jedno miejsce w lewo, a odjęcie
$1$ wyrzuca cyfrę, która przeszła przez przecinek, więc $\num{0.1011}$ staje
się $\num{0.011}$, potem $\num{0.11}$, potem $\num{0.1}$. Odwzorowanie
odczytuje swój start po jednej cyfrze dwójkowej naraz.

\begin{think}
Policz odwzorowanie podwajające ręcznie od $\tfrac{1}{10}$, na ułamkach,
przez pięć kroków. Co jego orbita robi potem na zawsze?
\end{think}
\reveal{$\tfrac{1}{5}$, $\tfrac{2}{5}$, $\tfrac{4}{5}$, $\tfrac{3}{5}$,
$\tfrac{1}{5}$: krąży po tych samych czterech wartościach bez końca i nigdy
nie dochodzi do $0$.}
Tymczasem komputerowe $\num{0.1}$ to liczba całkowita podzielona przez
$2^{\val{ch16.tenth.power}}$: ma ostatnią cyfrę dwójkową, a każdy krok
wypycha jedną cyfrę za przecinek.

\begin{think}
Jaka będzie komputerowa orbita $\num{0.1}$ po sześćdziesięciu krokach?
\end{think}
\reveal{Dokładnie $0$. Gdy ostatnia cyfra dwójkowa odejdzie, w kroku
\val{ch16.collapse.tenth}, nic nie zostaje, a $0$ podwojone to $0$ na
zawsze.}

\pyregion{code/ch16/doubling.py}{race}{Odwzorowanie podwajające od jednej dziesiątej: w komputerze i na dokładnych ułamkach}{lst:ch16-doubling}

\transcript{ch16-doubling}

\noindent
Przez czterdzieści kroków obie orbity zgadzają się do sześciu miejsc po
przecinku; potem komputer odpływa i ginie. Obliczona orbita może się mylić
w całym swoim charakterze: pętla zamieniła się w upadek do zera.

Najdziwniejsze jest to, że arytmetyka komputera nie popełniła żadnego błędu.
Podwojenie liczby podwójnej precyzji i odjęcie $1$ nigdy nie wymagają
zaokrąglenia, a skrypt pomiarowy rozdziału sprawdza krok po kroku, że orbita
komputera jest dokładnie prawdziwą orbitą liczby, którą zapamiętał. Błędny
jest start. Liczba, którą da się zapamiętać, ma skończenie wiele cyfr
dwójkowych, więc każdy taki start zostaje pochłonięty przez $0$, a prawie
każdy prawdziwy start ma nieskończenie wiele cyfr i nie zostaje pochłonięty
nigdy. Z każdego
z \val{ch16.collapse.starts} startów $\tfrac{1}{997}$, $\tfrac{2}{997}$ i tak
dalej orbita komputera zginęła po liczbie kroków od \val{ch16.collapse.min}
do \val{ch16.collapse.max}, najczęściej po \val{ch16.collapse.mode}.

\plotfig{ch16-collapse}{Odwzorowanie podwajające od jednej dziesiątej.
Prawda (puste kółka) krąży w pętli bez końca; komputer (kropki) zgadza się z
nią przez mniej więcej czterdzieści kroków, odpływa, a w kroku
\val{ch16.collapse.tenth} spada do $0$ na dobre.}{odwzorowanie podwajające
ginące w komputerze}

\begin{lab}{l16_01_collapse}{Kiedy ginie odwzorowanie podwajające?}
Napisz odwzorowanie podwajające, funkcję, która liczy kroki, aż orbita
będzie dokładnie $0$, i funkcję, która przewiduje tę liczbę bez uruchamiania
czegokolwiek, na podstawie \code{x.as\_integer\_ratio()}, które podaje
zapamiętaną liczbę jako liczbę całkowitą przez potęgę dwójki. Test
sprawdza, że obie zawsze się zgadzają.
\end{lab}

\section{Jedna reguła, dwa programy}
\label{sec:ch16-formulas}
\index{odwzorowanie logistyczne!w komputerze}\index{powtarzalność}

\noindent
Odwzorowanie podwajające nigdy nie zaokrągla. Odwzorowanie logistyczne z
rozdziału~\ref{ch:logistic} jest bardziej typowe: zaokrągla prawie każdy
krok. Oto ono przy $r = 4$, zapisane na dwa sposoby.

\pyregion{code/ch16/two_formulas.py}{rules}{Odwzorowanie logistyczne przy r = 4 zapisane na dwa sposoby}{lst:ch16-rules}

\noindent
Wymnóż nawias w pierwszym, a dostaniesz drugie:
$4x(1 - x) = 4x - 4x^{2}$. Na papierze to jedna reguła.

\begin{think}
Uruchom oba od $\num{0.1}$ na tym samym komputerze i pozwól im działać przez
siedemdziesiąt kroków. Czy oba programy wydrukują te same liczby?
\end{think}
\reveal{Nie. Różnią się na ostatniej cyfrze dwójkowej od kroku
\val{ch16.formulas.first}, a w kroku \val{ch16.formulas.apart} dzieli je już
więcej niż $\num{0.1}$.}
Pierwszy wzór zaokrągla $1 - x$, a potem iloczyn; drugi zaokrągla $4x^{2}$,
a potem różnicę. Każdy błąd to mniej więcej jedna część na $10^{16}$, a
odwzorowanie podwaja go w każdym kroku. Prawa kolumna poniżej to prawda:
orbita dokładnie jednej dziesiątej, liczona ze stu cyframi znaczącymi.

\transcript{ch16-two-formulas}

\noindent
Który wzór ma rację? Oba, przez czterdzieści kroków. W kroku
\val{ch16.wrong.a} pierwszy odszedł od prawdy, a w kroku \val{ch16.wrong.b}
także drugi. Żaden wzór nie jest błędny; każdy jest poprawny przez mniej
więcej pięćdziesiąt kroków, a potem podąża już za inną orbitą.

\begin{trapbox}
\textbf{\enquote{Program deterministyczny zawsze daje ten sam wynik.}}
Uruchom ten sam program na tej samej maszynie, a da, co do ostatniej cyfry.
Ale zmień komputer, kompilator, bibliotekę matematyczną albo kolejność
dodawania długiej sumy, a zmienisz to, które zaokrąglenia zachodzą \dash{}
dokładnie tak jak dwa wzory powyżej. Kompilatory mogą dla szybkości
przestawiać działania albo łączyć mnożenie i dodawanie w jeden krok z jednym
zaokrągleniem. W większości programów różnica zostaje na ostatniej cyfrze. W
chaotycznym dociera do pierwszej.
\end{trapbox}

\noindent
Pięćdziesiąt to nie przypadek i da się je przewidzieć.

\begin{think}
Przy $r = 4$ błąd mniej więcej podwaja się w każdym kroku, co zmierzył
rozdział~\ref{ch:butterfly}. Błąd zaokrąglenia liczby podwójnej precyzji to
około $10^{-16}$, a dziesięć podwojeń mnoży błąd mniej więcej przez tysiąc.
Po mniej więcej ilu krokach błąd $10^{-16}$ urośnie do $\num{0.1}$?
\end{think}
\reveal{Po około pięćdziesięciu. Od $10^{-16}$ do $10^{-1}$ to piętnaście
czynników dziesięć, a każde trzy czynniki dziesięć to mniej więcej dziesięć
podwojeń. Horyzont z rozdziału~\ref{ch:lyapunov}, z błędem zaokrąglenia
jako błędem startu, daje \val{ch16.predict.double}.}
Jednym zdaniem: przy $r = 4$ odwzorowanie zużywa jedną cyfrę dwójkową
dokładności na krok, więc \val{ch16.bits.double} cyfry liczby podwójnej
precyzji starczają na mniej więcej pięćdziesiąt kroków. \emph{Pojedyncza
precyzja}, której używają karty graficzne i wiele programów uczenia
maszynowego, bo zajmuje połowę pamięci, zachowuje \val{ch16.bits.single} cyfry
dwójkowe, czyli około \val{ch16.decimal.single} dziesiętnych; ta sama arytmetyka przewiduje \val{ch16.predict.single} kroków, a
wytrzymuje ona \val{ch16.wrong.single}.

\plotfig{ch16-parting}{Odległość od prawdziwej orbity jednej dziesiątej, w
skali logarytmicznej, dla odwzorowania logistycznego przy $r = 4$. Błąd
każdego przebiegu podwaja się co krok, wzdłuż kropkowanej linii, aż wypełni
cały przedział: po mniej więcej dwudziestu krokach w pojedynczej precyzji,
pięćdziesięciu w podwójnej.}{dwie precyzje, które się rozchodzą}

\begin{lab}{l16_02_single}{Pojedyncza przeciw podwójnej}
Zaokrąglij liczbę do pojedynczej precyzji, pakując ją do czterech bajtów i
odczytując z powrotem (plik startowy pokazuje, jak). Napisz odwzorowanie
logistyczne przy $r = 4$ z każdym działaniem zaokrąglanym w ten sposób,
uruchom je obok wersji podwójnej precyzji i znajdź pierwszy krok, w którym
różnią się o więcej niż $\num{0.1}$.
\end{lab}

\begin{worldbox}
Uczenie maszynowe spotyka to codziennie. Sieć neuronową trenuje się,
dodając miliony małych liczb na karcie graficznej, której równoległe wątki
mogą za każdym razem kończyć w innej kolejności, więc suma zaokrągla się
inaczej przy każdym przebiegu. Dwa przebiegi tego samego programu z tym samym
ziarnem losowym mogą dać sieci, których wagi różnią się wszędzie \dash{} i
które zwykle wypadają w testach mniej więcej tak samo. Dokumentacja
PyTorcha, szeroko używanej biblioteki, mówi wprost, że w pełni powtarzalne
wyniki nie są gwarantowane między platformami i wersjami, i oferuje
przełącznik wymuszający stałą kolejność działań, kosztem szybkości.
\end{worldbox}

\section{Dlaczego więcej cyfr cię nie uratuje}
\label{sec:ch16-exact}
\index{arytmetyka dokładna}

\noindent
Oczywiste lekarstwo to przestać zaokrąglać. Moduł \pkg{fractions} Pythona
przechowuje każdą liczbę jako dokładny ułamek. Jedna dziesiąta zostaje
jedną dziesiątą, a $4x(1 - x)$ daje z niej $\tfrac{9}{25}$, potem
$\tfrac{576}{625}$, bez żadnego błędu.

\begin{think}
Od $\tfrac{9}{25}$ do $\tfrac{576}{625}$ mianownik zmienił się z $25$ na
$625 = 25 \times 25$. Jeśli jest podnoszony do kwadratu w każdym kroku, co
dzieje się z liczbą jego cyfr?
\end{think}
\reveal{W każdym kroku mniej więcej się podwaja: podniesienie liczby do
kwadratu mniej więcej podwaja liczbę jej cyfr.}

\pyregion{code/ch16/more_digits.py}{exact}{Dokładne ułamki: mianownik podnoszony do kwadratu w każdym kroku}{lst:ch16-exact}

\transcript{ch16-more-digits}

\noindent
Po dziesięciu krokach mianownik ma \val{ch16.exact.ten} cyfr. Po
pięćdziesięciu miałby ich około \val{ch16.exact.fifty}, a samo jego
zapisanie zajęłoby jakieś \val{ch16.exact.fifty.tb} terabajtów. Arytmetyka
dokładna liczy każdy krok bez błędu i nie może skończyć. Ostatnie cztery
wiersze próbują lekarstwa praktycznego: trzymać stałą, ale większą liczbę
cyfr, modułem \pkg{decimal} Pythona, i porównać z przebiegiem liczonym z
tysiącem cyfr. Szesnaście cyfr wystarcza na \val{ch16.right.sixteen} kroki,
jak w liczbie podwójnej precyzji; dwa razy więcej cyfr wystarcza na dwa razy
dłużej. Każda cyfra dziesiętna kupuje mniej więcej \val{ch16.right.per.digit}
kroku, bo jest warta mniej więcej \val{ch16.right.per.digit} cyfry dwójkowej.

\begin{trapbox}
\textbf{\enquote{Więcej miejsc po przecinku czyni symulację dokładną.}}
Czyni ją poprawną dłużej, o stałą liczbę kroków na cyfrę, nigdy na zawsze.
Żeby śledzić chaotyczną orbitę przez tysiąc kroków, trzeba nieść mniej więcej
trzysta cyfr dziesiętnych i znać start równie dokładnie, a tego nie da ci
żaden pomiar. To okrutna arytmetyka z rozdziału~\ref{ch:lyapunov} od strony
komputera: precyzja kupuje horyzont dłuższy o stałą, a nie proporcjonalnie.
\end{trapbox}

\section{Zła droga, dobry obraz}
\label{sec:ch16-shadow}
\index{cieniowanie}

\noindent
Każde długie chaotyczne obliczenie w tej książce podążało więc za złą
orbitą. Czy narysowało zły obraz? Sposób, by się przekonać, opiera się na
jednej obserwacji.

\begin{think}
Odwzorowanie logistyczne uruchomione do przodu podwaja mały błąd w każdym
kroku, przeciętnie. Co robi z małym błędem uruchomienie go wstecz, od każdej
wartości do poprzedniej?
\end{think}
\reveal{Przeciętnie go połowi. Co krok do przodu rozciąga, to krok wstecz
ściska.}
Zacznij więc od końca obliczonej orbity i idź wstecz. Do każdej wartości
prowadzą w jednym kroku dwie liczby, jedna poniżej $\tfrac{1}{2}$ i jedna
powyżej; bierz tę po tej samej stronie co wartość obliczona. Marsz wstecz
zmniejsza błędy, więc trzyma się blisko obliczonej orbity aż do startu, a to,
co znajduje, liczone z czterystoma cyframi, jest prawdziwą orbitą reguły.

\pyregion{code/ch16/shadow.py}{back}{Marsz wstecz po obliczonej orbicie w poszukiwaniu prawdziwej orbity, za którą podąża}{lst:ch16-back}

\transcript{ch16-shadow}

\noindent
Obliczona orbita $\num{0.1}$ odeszła od prawdy po \val{ch16.wrong.a}
krokach. Ale start różniący się od jednej dziesiątej o mniej więcej
\val{ch16.shadow.start} ma dokładną orbitę, która trzyma się w odległości
\val{ch16.shadow.gap} od każdej z \val{ch16.shadow.steps} obliczonych
wartości. Komputer nie podążał za orbitą, o którą go poproszono. Podążał,
wiernie, za orbitą startu tak bliskiego, że żaden pomiar nie odróżniłby ich
od siebie. Ta prawdziwa orbita to jego \emph{cień}.

\begin{rigourbox}
\emph{Lemat o cieniowaniu}, udowodniony przez Dmitrija Anosowa i Rufusa
Bowena, dotyczy układów, które w każdym punkcie rozciągają w niektórych
kierunkach i ściskają w innych, w tempie, które nigdy nie słabnie (układów
\emph{jednostajnie hiperbolicznych}, takich jak odwzorowanie podwajające).
Mówi, że jakkolwiek blisko chcesz, by obliczona orbita miała swój cień,
wystarczająco mały błąd na krok to gwarantuje, na zawsze. Odwzorowanie
logistyczne przy $r = 4$ nie całkiem jest takie: w pobliżu $x = \tfrac{1}{2}$
jego krzywa jest płaska i tam marsz wstecz rozciąga. Dla niego cieniowanie
potwierdzono starannymi obliczeniami na długich, ale skończonych odcinkach,
zwłaszcza Hammel, Yorke i Grebogi w 1987 roku. Ta książka nie dowodzi ani
jednego, ani drugiego; dowody są w podręcznikach układów dynamicznych dla
zaawansowanych.
\end{rigourbox}

\noindent
Teraz zapłata. Cień jest prawdziwą orbitą, więc spędza czas tak, jak
prawdziwe orbity. Mierzą to ostatnie dwa wiersze wydruku: przez milion
kroków pierwszy wzór spędza ułamek \val{ch16.stats.a} czasu poniżej
$\num{0.1}$, drugi \val{ch16.stats.b}, a dokładna odpowiedź, ze wzoru,
którego ta książka nie wyprowadza, to \val{ch16.stats.exact}. Dwa programy, których drogi rozeszły
się po pięćdziesięciu krokach, zgadzają się co do tego, gdzie orbita
mieszka. Dlatego ułamkom czasu na skrzydłach z rozdziału~\ref{ch:lorenz} i
statystykom klimatu z rozdziału~\ref{ch:weather} można ufać, choć żadnej
pojedynczej obliczonej drodze nie można.

\plotfig{ch16-statistics}{Z lewej: gdzie dwa przebiegi podwójnej precyzji
po milion kroków spędzają czas, wobec dokładnej odpowiedzi. Z prawej: udział
kroków poniżej \num{0.1} w miarę wydłużania przebiegów. Przebiegi podwójnej
precyzji ustalają się na wartości dokładnej; przebieg pojedynczej precyzji,
uwięziony w pętli, na błędnej.}{statystyki, które przeżywają złą drogę}

\mermaidfig{ch16-shadow}{Obliczona orbita trzyma się blisko prawdziwej
orbity z pobliskiego startu, a statystyki tamtej orbity to te, które pokazuje
obraz.}{obliczona orbita i jej cień}

\section{Maszyny, które muszą się powtarzać}
\label{sec:ch16-random}
\index{liczby losowe}\index{liniowy generator kongruencyjny}

\noindent
Najprostsza różnica ze wszystkich: między $0$ a $1$ jest \val{ch16.count.single} liczb pojedynczej precyzji i
\val{ch16.count.double} liczb podwójnej precyzji: ogromnie dużo, ale
skończenie wiele. Obliczona orbita musi więc prędzej czy później trafić na
wartość, którą już spotkała, a ponieważ reguła za każdym razem daje tę samą
następną wartość, od tej chwili krąży po tej samej pętli bez końca.

W pojedynczej precyzji pętle są dość krótkie, by je znaleźć. Od $\num{0.1}$
odwzorowanie logistyczne przy $r = 4$ błądzi przez \val{ch16.loop.tail}
kroki, a potem krąży bez końca po pętli \val{ch16.loop.length} wartości;
\val{ch16.loop.starts} różnych startów kończy w zaledwie
\val{ch16.loop.kinds} różnych pętlach. Pętla psuje też statystyki: przebieg
pojedynczej precyzji spędza ułamek \val{ch16.stats.single} czasu poniżej
$\num{0.1}$, a nie \val{ch16.stats.exact}, co pokazuje
rysunek~\ref{fig:ch16-statistics}. Liczby podwójnej precyzji też mają pętle,
ale o wiele za długie, by do nich dojść, i to jest praktyczny powód, by
chaos liczyć w podwójnej precyzji.

\begin{think}
Odwzorowanie logistyczne przy $r = 4$ nigdy się nie powtarza, pokrywa cały
przedział od $0$ do $1$ i na pierwszy rzut oka wygląda losowo. Czy byłoby
dobrym generatorem liczb losowych?
\end{think}
\reveal{Nie, z trzech powodów, które ten rozdział zmierzył. Jego wartości
gromadzą się przy $0$ i $1$: mniej więcej jedną piątą czasu spędza poniżej
$\num{0.1}$, a nie jedną dziesiątą. Każda wartość jest wyznaczona przez
poprzednią. A w komputerze wpada w pętlę.}

\begin{trapbox}
\textbf{\enquote{Chaotyczne odwzorowanie to dobry generator liczb losowych.}}
Chaos czyni przyszłość nieprzewidywalną dla kogoś, kto nie zna dokładnie
startu, a to część tego, czego potrzebuje losowość, nie całość. Dobry
generator musi też rozkładać wartości równomiernie, nie pokazywać wzoru
między kolejnymi wartościami i powtarzać się dopiero po pętli, której nikt
nigdy nie zdoła obejść. Te własności projektuje się za pomocą teorii liczb i
sprawdza testami statystycznymi; chaos nie daje ich za darmo.
\end{trapbox}

\noindent
Generator liczb losowych to deterministyczna reguła na skończonym zbiorze
wartości, zbudowana tak, by miała jedną ogromną pętlę. Najprostszy rodzaj,
\emph{liniowy generator kongruencyjny}, to jeden wiersz: pomnóż przez $a$,
dodaj $c$ i zachowaj resztę z dzielenia przez $m$. Dla $a = 5$, $c = 3$ i
$m = 16$, od $0$ dostajesz $3$, potem $5 \times 3 + 3 = 18$, czyli resztę
$2$, potem $13$, $4$, $7$, $6$, $1$, $8$, $11$, $10$, $5$, $12$, $15$, $14$,
$9$ i znowu $0$: wszystkie szesnaście możliwych wartości przed pierwszym
powtórzeniem. Odwzorowanie podwajające jest generatorem tego rodzaju z
$a = 2$ i $c = 0$, a na $k$ cyfrach dwójkowych ginie w $k$ krokach \dash{}
to upadek z podrozdziału~\ref{sec:ch16-doubling} w nowym przebraniu. Własny
generator Pythona, Mersenne Twister Makoto Matsumoto i Takujiego Nishimury
(1998), ma pętlę długości $2^{19937} - 1$. Ponieważ jest deterministyczny,
ta sama wartość początkowa, \emph{ziarno}, daje za każdym razem te same
\enquote{losowe} liczby i tak właśnie losowy eksperyment staje się
powtarzalny.

\begin{lab}{l16_03_loops}{Każda orbita kończy się pętlą}
Napisz liniowy generator kongruencyjny i funkcję, która dla dowolnej reguły
i startu znajduje, po ilu krokach orbita wchodzi w pętlę i jak długa jest ta
pętla. Wypróbuj ją na generatorze powyżej, a potem sam znajdź pętlę
odwzorowania logistycznego w pojedynczej precyzji.
\end{lab}

\noindent
Praktyczna lekcja mieści się w jednym zdaniu: nigdy nie ufaj pojedynczemu
długiemu przebiegowi chaotycznego obliczenia. Uruchom je jeszcze raz z
przestawionym wzorem, z większą i mniejszą precyzją, z mniejszym krokiem, i
wierz temu, co przetrwa. Droga dalej niż pięćdziesiąt kroków nie przetrwa.
Statystyki przetrwają.

\begin{summarybox}
\begin{itemize}
\item Liczba \textbf{podwójnej precyzji} zachowuje \val{ch16.bits.double}
  cyfry dwójkowe znaczące. Większość ułamków dziesiętnych, w tym
  $\num{0.1}$, zapisuje się niedokładnie, a każda suma i iloczyn są
  \textbf{zaokrąglane}.
\item \textbf{Odwzorowanie podwajające} odczytuje swój start po jednej
  cyfrze dwójkowej, więc w komputerze każda orbita ginie w $0$; prawdziwa
  orbita jednej dziesiątej krąży w pętli bez końca.
\item Dwa programy dla jednej reguły zaokrąglają inaczej i się rozchodzą.
  Przy $r = 4$ odwzorowanie logistyczne zużywa \textbf{jedną cyfrę dwójkową
  na krok}: mniej więcej pięćdziesiąt kroków w podwójnej precyzji,
  dwadzieścia w pojedynczej.
\item \textbf{Arytmetyka dokładna} podwaja liczbę cyfr w każdym kroku i nie
  może skończyć; \textbf{więcej cyfr} kupuje stałą liczbę kroków na cyfrę.
\item Obliczona orbita ma \textbf{cień}, prawdziwą orbitę z pobliskiego
  startu, więc jej \textbf{statystyki} są poprawne, choć droga jest błędna.
\item Skończona maszyna sprawia, że każda orbita kończy się \textbf{pętlą};
  generator liczb losowych to reguła zbudowana tak, by ta pętla była
  ogromna.
\end{itemize}
\end{summarybox}

\canyou

\begin{checkpoint}
\item Dlaczego w komputerze $\num{0.1} + \num{0.2}$ nie równa się
  $\num{0.3}$?
  \answerto{Ani $\num{0.1}$, ani $\num{0.2}$ nie ma dokładnej postaci
  dwójkowej, więc każda jest zapamiętana jako najbliższa liczba podwójnej
  precyzji, a suma jest jeszcze raz zaokrąglana.}
\item Komputer zapamiętuje $\num{0.3}$ jako liczbę całkowitą przez
  $2^{\val{ch16.collapse.three}}$. Kiedy jej orbita w odwzorowaniu
  podwajającym dochodzi do $0$ i co robi prawdziwa orbita $\tfrac{3}{10}$?
  \answerto{W kroku \val{ch16.collapse.three}, jeden krok na cyfrę dwójkową.
  Prawdziwa orbita idzie przez $\tfrac{3}{5}$, $\tfrac{1}{5}$,
  $\tfrac{2}{5}$, $\tfrac{4}{5}$, $\tfrac{3}{5}$ i krąży w tej pętli bez
  końca.}
\item Pojedyncza precyzja zachowuje \val{ch16.bits.single} cyfry dwójkowe.
  Przez mniej więcej ile kroków odwzorowania logistycznego przy $r = 4$
  jest dobra?
  \answerto{Przez około dwadzieścia (zmierzone: \val{ch16.wrong.single}):
  odwzorowanie podwaja błąd zaokrąglenia w każdym kroku, zużywając jedną
  cyfrę dwójkową na krok.}
\item Liczony dokładnie od $\tfrac{1}{10}$ mianownik ma po dziesięciu krokach
  \val{ch16.exact.ten} cyfr. Mniej więcej ile po dwudziestu?
  \answerto{Około \val{ch16.exact.twenty}: liczba cyfr podwaja się w każdym
  kroku, a dziesięć podwojeń mnoży ją mniej więcej przez tysiąc.}
\item Po milionie kroków chaotycznej symulacji czemu możesz ufać: ostatniej
  wartości czy ułamkowi czasu spędzonego poniżej $\num{0.1}$?
  \answerto{Ułamkowi czasu. Droga podąża za inną orbitą mniej więcej od
  pięćdziesiątego kroku, ale trzyma się blisko prawdziwej orbity z
  pobliskiego startu, a ta spędza czas tak, jak prawdziwe orbity.}
\item Dlaczego każda orbita liczona na maszynie musi się w końcu powtórzyć?
  \answerto{Maszyna mieści tylko skończenie wiele liczb, więc orbita musi
  wrócić do jakiejś wartości, a reguła powtarza wtedy wszystko, co po niej
  następuje.}
\end{checkpoint}
